\section{Introduction}

\label{sec:intro}

The rapid growth of AI workloads has created unprecedented demand for compute resources. Training a single frontier language model may require thousands of GPU-hours across multiple data centers, while inference serving demands low-latency access to heterogeneous accelerator types. Centralized cloud providers (AWS, GCP, Azure) dominate this market but impose rigid pricing, vendor lock-in, and geographic constraints. Decentralized compute marketplaces~\cite{golem2016,akash2020,rendernetwork2021} have emerged as alternatives, but face fundamental challenges in price discovery for resources that are inherently non-fungible and multi-dimensional.

\subsection{Limitations of Existing AMMs}

Automated market makers have revolutionized decentralized finance by enabling permissionless, oracle-free token exchange~\cite{adams2020uniswap,egorov2019stableswap,martinelli2019balancer}. The constant product market maker (CPMM), popularized by Uniswap~\cite{adams2021uniswapv3}, maintains the invariant $x \cdot y = k$ for a pair of token reserves $(x, y)$. While elegant, this formulation assumes:

\begin{enumerate}[leftmargin=1.4em]
    \item \textbf{Homogeneity}: Tokens within each reserve are fungible.
    \item \textbf{Two-dimensionality}: Each pool handles exactly two assets.
    \item \textbf{Time-independence}: The invariant function does not adapt to market conditions.
    \item \textbf{Memorylessness}: Each swap is independent; the AMM has no state beyond reserves.
\end{enumerate}

AI compute resources violate all four assumptions. GPU-hours from an A100 are not equivalent to those from an H100. A training job requires simultaneous allocation of GPU, memory, bandwidth, and storage. Market conditions fluctuate on timescales of minutes as training runs begin and end. And the quality of a compute provider, measured by uptime, latency, and attestation history, is a stateful quantity.

\subsection{Physics-Inspired Market Design}

Classical Hamiltonian mechanics~\cite{arnold1989mathematical,goldstein2002classical} provides a natural mathematical framework for dynamical systems with conservation laws. In Hamiltonian mechanics, the state of a system is described by generalized coordinates $\bm{q}$ and conjugate momenta $\bm{p}$, evolving according to Hamilton's equations:
\begin{equation}
\dot{q}_i = \frac{\partial \mathcal{H}}{\partial p_i}, \quad \dot{p}_i = -\frac{\partial \mathcal{H}}{\partial q_i}.
\label{eq:hamilton}
\end{equation}
The Hamiltonian $\mathcal{H}(\bm{q}, \bm{p})$ is conserved along trajectories, and the phase space flow preserves the symplectic two-form $\omega = \sum_i dp_i \wedge dq_i$.

We observe a deep structural analogy between Hamiltonian systems and market dynamics:
\begin{itemize}[leftmargin=1.1em]
    \item \textbf{Coordinates $\bm{q}$}: Resource reserves in the pool.
    \item \textbf{Momenta $\bm{p}$}: Demand pressures (credit reserves, order flow).
    \item \textbf{Hamiltonian $\mathcal{H}$}: Market invariant (replaces $xy = k$).
    \item \textbf{Conservation}: Prices emerge from the invariant without external oracles.
    \item \textbf{Symplecticity}: Volume preservation prevents information loss and ensures reversibility of the pricing function.
\end{itemize}

\subsection{Contributions}

This paper makes the following contributions:

\begin{enumerate}[leftmargin=1.4em]
    \item We formulate AMM pricing as a Hamiltonian system on a symplectic manifold and derive the complete equations of motion governing price evolution (\S\ref{sec:formulation}).
    \item We prove three conservation laws---energy, angular momentum, and phase space volume---and show how they map to economic invariants: no-arbitrage, balanced reserves, and information preservation (\S\ref{sec:conservation}).
    \item We construct symplectic integrators (Verlet, Ruth's 4th-order) that preserve the geometric structure under discrete swap events, ensuring long-term stability (\S\ref{sec:integrators}).
    \item We generalize to multi-asset pools with weighted resources and prove impermanent loss bounds strictly tighter than CPMM, Curve, and Balancer (\S\ref{sec:multiasset}).
    \item We introduce a MEV resistance mechanism based on non-commutative phase space flows (\S\ref{sec:mev}).
    \item We present simulation results over 180 days of synthetic data and a live testnet with 50 workers and 100 clients (\S\ref{sec:experiments}).
\end{enumerate}

